%----------------------------------------------------------------------------------------
%   moderncv — classic, widely-used, Overleaf built-in
%
%   编译方式：
%     Overleaf → 选 pdfLaTeX（moderncv 需要编译两次）
%     或者直接"New Project" → 搜 "moderncv" → 起模板 → 粘贴下面内容
%
%   moderncv.cls 在 Overleaf / TeX Live 中自带，不需要额外下载。
%----------------------------------------------------------------------------------------

\documentclass[11pt, a4paper]{moderncv}

% --- Choose a style:
%   classic  — traditional, no photo
%   banking  — header with coloured bar, very professional
%   casual   — relaxed, coloured header
%   oldstyle — serif, classic academic
\moderncvstyle{banking}

% --- Colour (optional themes: blue, orange, green, red, grey, black)
\definecolor{color0}{rgb}{0.17,0.24,0.31}   % deep navy — match homepage accent
\definecolor{color1}{rgb}{0.17,0.24,0.31}
\definecolor{color2}{rgb}{0.35,0.42,0.49}

% --- Page geometry
\usepackage[scale=0.85]{geometry}

%----------------------------------------------------------------------------------------
%   Personal data — no photo, no phone
%----------------------------------------------------------------------------------------
\name{Linsheng}{He}
\title{Ph.D.\ Candidate in Public Administration}
\address{School of Public Policy \& Management}{Tsinghua University}{Beijing, China}
\email{hels21@mails.tsinghua.edu.cn}
\homepage{linshenghe.github.io}
\social[github]{linshenghe}
\extrainfo{National Scholarship; ASPA-SCPA Best Student Paper Award}

\begin{document}

\makecvtitle

%----------------------------------------------------------------------------------------
%   Education
%----------------------------------------------------------------------------------------
\section{Education}

\cventry{2021.09 -- Expected 2026.12}
  {Ph.D.\ Candidate, Public Administration}
  {Tsinghua University}
  {Beijing, China}
  {}
  {
    \textbf{Advisor:} Prof.\ Yixin Dai (Director, Institute of Politics and Public Policy)\\
    \textbf{Dissertation:} Communication Mechanisms and Conflict Resolution in Collaborative Governance\\
    \textbf{Honors:} National Scholarship; Tsinghua Comprehensive Excellence Scholarship
  }

\cventry{2017.09 -- 2021.06}
  {B.A.\ in Public Affairs Management}
  {Beijing Normal University}
  {Beijing, China}
  {}
  {}

%----------------------------------------------------------------------------------------
%   Research Interests
%----------------------------------------------------------------------------------------
\section{Research Interests}

\cvitem{}{\begin{itemize}
  \item Policy design, policy evaluation, and the policy process
  \item Behavioral public management and experimental methods
  \item Government communication, public attitudes, and policy feedback
  \item Collaborative governance and cross-sector policy coordination
\end{itemize}}

%----------------------------------------------------------------------------------------
%   Publications
%----------------------------------------------------------------------------------------
\section{Publications}

\cvsubsection{English-language Journals}

\cvitem{2025}{
  \textbf{He, L.}, Dai, Y., \& Guo, Y.\ \textit{Aligning Behavioral Assumptions Underlying Policy Instruments.}
  \textit{Public Administration}, Online First.
  \hfill \textbf{SSCI, JCR Q1 / ABS 4}
}

\cvitem{2024}{
  Guo, Y., \& \textbf{He, L.\ (corresponding author)}.\ \textit{Communicating in Diverse Local Cultures.}
  \textit{Review of Policy Research}, 41(4), 654--678.
  \hfill \textbf{SSCI, JCR Q1}
}

\cvsubsection{Chinese-language Journals}

\cvitem{2021}{
  Guo, Y., Hong, J., \& \textbf{He, L.}\ \textit{A Policy Feedback Explanation of Government Procurement of Facial Recognition Technology.}
  \textit{Journal of Public Administration Review} [公共行政评论], 14(05), 159--177.
  \hfill \textbf{CSSCI}
}

\cvitem{2020}{
  Guo, Y., \textbf{He, L.}, \& Su, J.\ \textit{``Instruments--Narratives--Feedback'': A Behavioral Public Policy Research Framework.}
  \textit{Chinese Public Administration} [中国行政管理], (05), 71--78.
  \hfill \textbf{CSSCI}
}

\cvitem{2020}{
  Guo, Y., Hong, J., \& \textbf{He, L.}\ \textit{The Policy Intervention Effect on Public Service Motivation.}
  \textit{Leadership Science Forum} [领导科学论坛], (23), 20--30.
}

%----------------------------------------------------------------------------------------
%   Working Papers
%----------------------------------------------------------------------------------------
\section{Working Papers}

\cvitem{R\&R, 2nd Round}{
  Hong, J., \& \textbf{He, L.\ (corresponding author)}.\
  ``Not All Ties Are Equal: Actor- and Type-Specific Social Capital in Coproduction.''
  \textit{Public Management Review}.
}

\cvitem{Under Review}{
  \textbf{He, L.}, \& Dai, Y.\
  ``The Micro-foundations of Reputational Persistence.''
}

\cvitem{Under Review}{
  Chen, Z., \& \textbf{He, L.\ (corresponding author)}.\
  ``Anti-Corruption, Trust, and Compliance in Public Health Crises.''
  \textit{Humanities and Social Sciences Communications}.
}

\cvitem{In Preparation}{
  Fu, X., \textbf{He, L.}, Zhong, W., \& Dai, Y.\
  ``Public Risk Perception during the COVID-19 Pandemic: A Meta-analysis.''
}

%----------------------------------------------------------------------------------------
%   Conference Presentations
%----------------------------------------------------------------------------------------
\section{Selected Conference Presentations}

\cventry{2026}{ASPA Annual Conference (Hollywood); IRSPM Conference (Perth)}{}{}{}{}

\cventry{2025}{PMRC (Seoul); ICCPS (Beijing)}{}{}{}{}

\cventry{2024}
  {PMRC (Seattle); ICCPS (Beijing) --- Best Paper Award}
  {}
  {}
  {}
  {}

\cventry{2023}{AP-PPN Conference (Hong Kong)}{}{}{}{}

\cventry{2022}{AAPA Annual Conference (Shanghai) --- 2 papers}{}{}{}{}

\cventry{2021}{APPAM Fall Research Conference (Online) --- Panel Organizer}{}{}{}{}

\cventry{2020}{IPMN Conference (Online)}{}{}{}{}

\cventry{2019}{AP-PPN Conference}{}{}{}{}

%----------------------------------------------------------------------------------------
%   Teaching Experience
%----------------------------------------------------------------------------------------
\section{Teaching Experience}

\cventry{2022 -- 2026}
  {Teaching Assistant}
  {School of Public Policy \& Management, Tsinghua University}
  {}
  {}
  {
    9 courses including:
    \begin{itemize}
      \item Public Policy Analysis (English, IMPA core course)
      \item History of Public Administration Thought (Ph.D.\ core course)
      \item Frontiers in Risk Management; Advanced Academic Writing; Urban Development
    \end{itemize}
  }

%----------------------------------------------------------------------------------------
%   Honours \& Awards
%----------------------------------------------------------------------------------------
\section{Honours \& Awards}

\cvitem{2026}{ASPA-SCPA \textbf{Best Student Paper Award}}
\cvitem{2025}{National Scholarship}
\cvitem{2024}{ICCPS \textbf{Best Paper Award}}
\cvitem{2024}{Tsinghua University Comprehensive Excellence Scholarship (Second Class)}
\cvitem{2022}{Tsinghua University Comprehensive Excellence Scholarship (First Class)}
\cvitem{2021}{\textbf{Best Paper Award}, 5th Young Scholars Forum, \textit{Journal of Public Administration Review}}

%----------------------------------------------------------------------------------------
%   Academic Service
%----------------------------------------------------------------------------------------
\section{Academic Service}

\cvitem{Editorial Assistant}{\textit{Public Administration}, 2024--present.}

\cvitem{Journal Reviewer}{\textit{Governance}; \textit{Public Management Review}; \textit{Information Polity}; \textit{Journal of Public Administration Review}.}

\cvitem{Memberships}{American Society for Public Administration (ASPA); International Research Society for Public Management (IRSPM).}

\end{document}
